\documentclass[11pt]{article}

\usepackage{amsmath,amssymb,amsthm}

\title{Cyclone: Terminal Rigidity Inequality}
\date{}

\newtheorem{theorem}{Theorem}

\begin{document}

\maketitle

\section*{Statement}

Let $G$ be a bounded–degree graph with maximum degree $\Delta$.

Let

\[
\mathrm{COR}(G)
\]

denote the cycle–overlap rank.

Let

\[
\operatorname{tp}_{FO^k}(B_R(v))
\]

denote the $\mathrm{FO}^k$ type of the radius-$R$ neighborhood of $v$.

\begin{theorem}[Cyclone]
For fixed parameters $k,\Delta,R$ there exists a constant
\[
C(k,\Delta,R)
\]
such that any graph satisfying

\[
\forall u,v \in V(G):
\operatorname{tp}_{FO^k}(B_R(u))=
\operatorname{tp}_{FO^k}(B_R(v))
\]

must satisfy

\[
\mathrm{COR}(G)\le C(k,\Delta,R).
\]

\end{theorem}

\section*{Interpretation}

The inequality asserts that FO$^k$ local homogeneity prevents
unbounded cycle–overlap complexity.

This forms the final rigidity obstruction used in the
Oblivion Atom and Chronos frameworks.

\end{document}
